\section{贯穿项目：做一个可解释的 C++20 本地推理服务}

\subsection{本册最终要让读者做成什么}

如果读完第三册只知道一些名词，例如 Transformer、KV cache、量化、GEMM、GPU、AI 编译器，那还不够。本册的贯穿项目要更具体：做一个受限但严肃的 C++20 本地推理服务原型。它不承诺一开始达到成熟框架的吞吐，也不承诺支持所有模型格式；它要做到的是把模型资产、编译检查、推理运行时、后端优化、正确性 oracle 和性能报告接成一条可解释链路。

\begin{keyidea}
贯穿项目的目标不是“能生成几句话”，而是“每个 token 为什么这样算、用了多少内存、读了多少权重、KV cache 怎样增长、量化误差从哪里来、性能瓶颈卡在哪一层”都能被追溯。
\end{keyidea}

项目完成后的最小能力如下：

\begin{lstlisting}[numbers=none,caption={第三册贯穿项目的最终能力}]
load:
  read model manifest
  verify tensor shape, dtype, quantization and tokenizer contract
  build typed graph
  lower into CPU/GPU executable plan

run:
  tokenize prompt
  run prefill
  allocate and append KV cache
  decode token by token
  stream token events
  stop, cancel or reject by resource policy

prove:
  compare reference logits
  compare optimized operator outputs
  report quantization error
  report TTFT, decode latency, memory peak and fallback
\end{lstlisting}

这个项目能支撑真实业务原型，例如本地问答、RAG 检索增强回答、批量摘要、代码补全或内部文档助手。但业务能力必须建立在底层合同之上：上下文长度要能预算，取消要能释放 KV cache，trace 要能定位慢点，错误要能说明是模型资产、后端能力、量化 profile 还是请求参数的问题。

\subsection{不是完整框架，但必须完整闭环}

一个学习项目不需要一开始覆盖所有工业功能。可以暂缓的东西包括：全部模型方言、分布式推理、复杂动态 batching、PagedAttention 的生产级调度、所有 GPU kernel、所有低比特格式、服务多租户隔离。不能暂缓的是闭环本身。

\begin{center}
\begin{tabular}{p{0.22\linewidth}p{0.33\linewidth}p{0.31\linewidth}}
\toprule
层次 & 可以先做受限版本 & 不能缺失的合同 \\
\midrule
模型导入 & 只支持一种模型切片或一种 manifest & shape、dtype、量化、tokenizer 必须验证 \\
推理数学 & 先跑小模型或一层切片 & reference logits 必须可复现 \\
CPU 后端 & 先标量，再 SIMD，再线程 & operator oracle 和 profiler 必须存在 \\
GPU 后端 & 先建立 device/stream/fallback 合同 & copy、launch、sync 账本必须可见 \\
服务接口 & 先单进程 HTTP 或 CLI & admission、取消、trace、finish reason 必须明确 \\
\bottomrule
\end{tabular}
\end{center}

闭环的含义是：输入资产先被验证，运行过程能被 trace，输出结果能被 oracle 比较，性能数字能回到具体阶段。若只写一个 \code{generate(prompt)}，即使它能输出文本，也无法回答“为什么慢”“为什么漂移”“为什么显存涨了”“为什么换量化格式后质量变差”。

\subsection{可执行交付物}

贯穿项目最终应该形成四类交付物。

\topic{命令行工具。} 命令行最适合教学和回归。它不需要复杂服务依赖，容易固定输入和输出。

\begin{lstlisting}[numbers=none,caption={贯穿项目命令行工具}]
ai-inspect-model --model ./model
ai-compile-plan --model ./model --backend cpu --quant q4_group64
ai-generate --model ./model --prompt prompts/basic.txt --max-new-tokens 64
ai-bench --model ./model --prompt-set prompts/bench.jsonl
ai-oracle --fixture fixtures/tiny_decoder.json
\end{lstlisting}

\code{ai-inspect-model} 只读模型资产并输出 manifest 诊断；\code{ai-compile-plan} 只做 verifier、IR 和 lowering；\code{ai-generate} 跑端到端请求；\code{ai-bench} 输出性能报告；\code{ai-oracle} 对比 reference。把这些命令拆开，读者才能知道问题发生在哪一层。

\topic{本地服务。} 服务形态可以先是单进程 HTTP 或本地 RPC。它不追求复杂微服务体系，但请求和响应字段必须严肃。

\begin{lstlisting}[numbers=none,caption={本地推理服务接口}]
POST /v1/generate
  model_id
  prompt
  max_new_tokens
  max_context_tokens
  temperature
  top_k
  top_p
  backend_policy
  quant_profile_id
  stream
  trace_level

stream event:
  token | final | error | trace
\end{lstlisting}

业务层需要的不是一句“模型生成失败”。它要知道是上下文超限、队列满、内存预算不足、后端不可用、量化 profile 不支持，还是运行时被取消。服务接口因此必须和 admission control、session、KV cache、trace 和错误诊断绑定。

\topic{报告文件。} 每次运行都应该能输出机器可读报告。报告让优化可以被比较，而不是靠主观体感。

\begin{lstlisting}[numbers=none,caption={贯穿项目报告文件}]
reports/
  manifest-report.json
  plan-report.json
  oracle-report.json
  trace.jsonl
  benchmark-report.json
  release-report.json
\end{lstlisting}

\code{manifest-report} 说明模型资产是否可信；\code{plan-report} 说明每个 segment 选择的 backend、layout 和 workspace；\code{oracle-report} 说明误差；\code{trace} 记录运行事件；\code{benchmark-report} 记录性能；\code{release-report} 记录当前版本支持范围和已知限制。

\topic{测试夹具。} 没有固定 fixture，就无法形成教学和回归。夹具应从小到大：

\begin{lstlisting}[numbers=none,caption={测试夹具分层}]
fixtures/
  tokenizer_tiny/
  tensors_layout/
  quant_bytes/
  tiny_rmsnorm/
  tiny_linear/
  tiny_rope_attention/
  tiny_decoder_block/
  tiny_end_to_end/
  7b_shape_only/
\end{lstlisting}

\code{tiny\_end\_to\_end} 可以只包含极小参数和固定 token，不需要真实 7B 权重；\code{7b\_shape\_only} 用来验证真实规模 shape、KV 容量和 plan，不必包含完整权重。这样项目既能快速测试数学正确性，也能训练读者面对真实 7B 规模的资源账本。

\subsection{建议目录结构}

C++20 项目不能把所有东西塞进一个大文件。下面是一套随章节增长的目录结构。

\begin{lstlisting}[numbers=none,caption={贯穿项目目录结构}]
include/ai/
  model/        manifest, tensor index, tokenizer contract
  graph/        tensor desc, graph node, verifier facts
  quant/        quant desc, converter, calibration record
  compile/      pass pipeline, memory plan, lowering
  runtime/      session, KV cache, arena, sampler, scheduler
  backend/
    cpu/        capability, packed weight, kernels, thread plan
    gpu/        device buffer, stream, event, workspace, fallback
  oracle/       reference, comparison, error report
  report/       trace schema, benchmark schema, release schema

src/
  model/
  graph/
  quant/
  compile/
  runtime/
  backend/cpu/
  backend/gpu/
  oracle/
  report/

tools/
  inspect_model/
  compile_plan/
  generate/
  bench/
  oracle/

tests/
  model/
  graph/
  quant/
  runtime/
  backend/cpu/
  backend/gpu/
  oracle/
\end{lstlisting}

这些目录对应的是责任边界。模型导入不选择 SIMD kernel；后端 kernel 不解析 tokenizer；runtime 不重新做 graph verifier；oracle 不复用被测后端的可疑路径；report 只描述事实，不偷偷改变执行。目录边界清楚，后面扩 GPU、量化、服务接口时才不会互相污染。

\subsection{核心数据结构合同}

项目的核心不是类名，而是这些对象之间的合同。下面的伪代码只表达字段，不要求读者照抄实现。

\begin{lstlisting}[numbers=none,caption={模型与图的核心合同}]
TensorDesc:
  name
  role
  shape
  dtype
  layout
  quant_desc_optional
  file_range

OperatorDesc:
  id
  op_kind
  inputs
  outputs
  attributes
  state_effects

ModelManifest:
  model_id
  tokenizer_contract
  tensor_table
  layer_count
  hidden_size
  head_count
  kv_head_count
  rope_config
\end{lstlisting}

\code{TensorDesc} 连接文件、shape、dtype、layout 和量化；\code{OperatorDesc} 表达算子语义与状态读写；\code{ModelManifest} 记录模型级合同。后端只应该接收已经验证过的 descriptor。

\begin{lstlisting}[numbers=none,caption={运行时与后端的核心合同}]
ExecutableSegment:
  id
  op_or_fused_region
  backend
  input_values
  output_values
  state_effects
  kernel_desc
  workspace_bytes
  profiler_label

SessionState:
  request_id
  token_buffer
  model_position
  kv_cache
  activation_arena
  sampler_state
  trace

BackendResult:
  status
  output_values
  profiler_events
  fallback_reason_optional
\end{lstlisting}

\code{ExecutableSegment} 是编译器和后端的接口；\code{SessionState} 是服务请求和 runtime 的接口；\code{BackendResult} 是后端和 profiler/oracle 的接口。若这些对象缺失，工程就会靠隐式全局状态和临时约定运行，后面很难定位错误。

\subsection{一条请求的业务级状态机}

业务开发必须把生成过程做成状态机，而不是阻塞函数。

\begin{lstlisting}[numbers=none,caption={推理请求状态机}]
received
  -> validated
  -> admitted
  -> tokenized
  -> prefill_running
  -> decoding
  -> finishing
  -> completed

failure edges:
  rejected
  cancelled
  timed_out
  backend_failed
  oracle_failed_debug_only
\end{lstlisting}

每条边都有资源语义。进入 \code{admitted} 时要预留 KV cache 和 workspace；进入 \code{prefill\_running} 后要保证 plan 已准备好；进入 \code{decoding} 后每个 step 会追加 token 和 K/V；进入 \code{cancelled} 或 \code{timed\_out} 必须释放 session 资源；进入 \code{completed} 要写 usage 和 trace summary。

\begin{center}
\begin{tabular}{p{0.18\linewidth}p{0.38\linewidth}p{0.30\linewidth}}
\toprule
状态 & 持有资源 & 必须记录的事实 \\
\midrule
admitted & KV 预算、队列 slot & 接受原因、估计内存、计划 id \\
prefill & activation arena、workspace & prompt tokens、TTFT 分解 \\
decoding & KV cache、sampler state & step latency、context length、finish reason \\
finishing & trace buffer、usage & token 数、峰值内存、fallback \\
cancelled & 待释放资源 & 取消来源、已释放字节数 \\
\bottomrule
\end{tabular}
\end{center}

这张表直接回答“能不能做实例业务开发”。能，但前提是业务服务不是只调用一个黑箱函数，而是理解请求状态、资源预算、取消、错误和 trace。

\subsection{第一个可运行版本的验收}

第一个版本可以很小，但验收要严。

\begin{rubricbox}
\begin{itemize}
  \item 给定一个 tiny manifest，\code{ai-inspect-model} 能输出 tensor 表、shape 和 tokenizer 合同。
  \item 故意改坏 embedding 行数，verifier 能报出 tokenizer vocab 与 embedding shape 不匹配。
  \item 给定 tiny decoder fixture，reference 能输出固定 logits。
  \item \code{ai-generate} 能跑 prefill 和至少 8 个 decode step，并输出 token event。
  \item 超过 context limit 的请求被 admission control 拒绝，错误里包含实际 token 数和允许上限。
  \item 打开 trace 后，报告里能看到 tokenize、prefill、decode、sample 的耗时。
  \item 取消请求后，KV cache 预留字节数回到请求前。
\end{itemize}
\end{rubricbox}

如果这些验收都通过，项目已经不是玩具打印程序。它具备业务服务最小骨架，也具备继续扩 CPU SIMD、量化、GPU 和框架对标的证据基础。

\subsection{贯穿项目和后续内容怎样对应}

后续内容不是散讲。每一节都落到贯穿项目的某个对象。

\begin{center}
\begin{tabular}{p{0.24\linewidth}p{0.31\linewidth}p{0.29\linewidth}}
\toprule
主题 & 项目对象 & 验收证据 \\
\midrule
张量布局 & \code{TensorView}、\code{TensorDesc} & 地址公式、stride 测试 \\
GEMM/SIMD & \code{PackedWeight}、\code{MicroKernelDesc} & operator oracle、latency \\
量化 & \code{QuantDesc}、\code{QuantConverter} & byte oracle、误差报告 \\
KV cache & \code{KVCache}、\code{KVCacheDesc} & position 测试、容量报告 \\
AI 编译器 & \code{GraphNode}、\code{ExecutableSegment} & verifier、plan report \\
CPU/GPU 后端 & \code{BackendResult}、\code{StepTrace} & profiler、fallback report \\
工程验收 & \code{ReleaseReport}、\code{PerformanceGate} & CI、对标、回归门禁 \\
\bottomrule
\end{tabular}
\end{center}

读者读到任何小节，都应该能问三个问题：这个对象在项目里叫什么，谁生产它，谁消费它，怎样测试它。回答不了这三个问题的内容，不应该进入贯穿项目主线。

\subsection{本节验收线}

读完本节，应该明确第三册不是泛泛介绍 AI，而是围绕一个可解释本地推理服务推进：

\begin{itemize}
  \item 最终交付物包括 CLI、本地服务、报告文件和测试夹具。
  \item 项目边界包括模型导入、typed graph、量化、runtime、CPU/GPU 后端、oracle 和 profiler。
  \item 业务接口必须包含 admission、取消、finish reason、usage 和 trace。
  \item 最小版本可以受限，但 manifest、reference、KV cache、trace、错误诊断和资源释放不能缺失。
  \item 后续内容所有概念都要落到具体对象、生产者、消费者和验收证据。
\end{itemize}
